% !PRINTED 441
\weberpage{441}
%<<CONT
$\dfrac{n^2-1}{8}$ eine gerade, in den beiden letzten eine ungerade Zahl,
und demnach erhalten wir die genaue Formel (1):

\weq{3}{2^{\frac{n-1}{2}} \overset{\nu}{\Pi} \cos \frac{2\nu m \pi}{n} = (-1)^{\frac{n^2-1}{8}}.}

In der Formel (2) hängt das Vorzeichen von $m$ ab; es ist
positiv oder negativ, je nachdem unter den Zahlen $(\varrho)$ eine
gerade oder eine ungerade Anzahl negativer vorkommt.

Wir setzen nun, um die Abhängigkeit des Vorzeichens von
$m$ und $n$ in der Bezeichnung auszudrücken, für (2)

\weq{4}{2^{\frac{n-1}{2}} \overset{\nu}{\underset{1,\,\frac{n-1}{2}}{\Pi}} \sin \frac{2\nu m \pi}{n} = \left(\frac{m}{n}\right)\sqrt{n},}

worin

\[
\left(\frac{m}{n}\right) = \pm 1
\]

ist.

Das Symbol $\left(\dfrac{m}{n}\right)$ ist von einer anderen Seite her und in
speciellerer Fassung von \gesperrt{Legendre} in die Zahlentheorie eingeführt
 und von \gesperrt{Jacobi} verallgemeinert worden. Wir wollen
es das \gesperrt{Legendre'sche Symbol} nennen. Seine Bedeutung für
die Zahlentheorie wird sich gleich ergeben; zunächst haben
wir den Satz:

1. \gesperrt{Es ist} $\left(\dfrac{m}{n}\right)$ \gesperrt{gleich} $+1$ \gesperrt{oder gleich} $-1$,
\gesperrt{je nachdem die Anzahl der negativen unter den absolut
kleinsten Resten von} $m\nu$ \gesperrt{eine gerade oder eine ungerade
ist.}

Für $m = 1$ gilt, wie wir schon gesehen haben, das positive
Zeichen in der Formel (2); also haben wir, wie auch die
Formel (4) direct zeigt,

\[
2.\quad \left(\frac{1}{n}\right) = +1.
\]

Verwandeln wir $m$ in $-m$, so ändert sich in allen Factoren
des Productes auf der linken Seite von (4) das Vorzeichen; da
die Anzahl der Factoren $\dfrac{n-1}{2}$ beträgt, so ergiebt sich

\[
3.\quad \left(\frac{-m}{n}\right) = (-1)^{\frac{n-1}{2}} \left(\frac{m}{n}\right)
\]

\gesperrt{und für} $m = 1$
%CONT>>
